\section{Wirtinger words and the Alexander polynomial}
\label{sec:alex}
A meridian label stays fixed along an overpass. On an underpass a label $x$ becomes $S_y(x)$ on crossing the normal of the oriented overpass, whose label is $y$. Choose the normal so that a positive downward crossing with the left strand over gives bottom colors $(S_x(y),x)$. Reading backwards uses $S_y^{-1}$. In this section we let the labels themselves be formal generators. We obtain invariants even where the sets they present are infinite, by recording finite substitutions and finite verifications of the relations. The classical background is the Wirtinger presentation and the first elementary ideal \cite{alexander,hiller,colorknots}.

\begin{lemma}[Eliminating local labels]\label{lem:eliminate}
Suppose $S_y$ is invertible, fixes $y$, and $S_z S_y=S_{S_z(y)}S_z$. If the inverse and these identities are witnessed by formal computations in the labels, either side of a Reidemeister replacement gives the same relations on the outside labels after explicitly eliminating the inside labels.
\end{lemma}
\begin{proof}
A variable appearing as $v=W$ with $W$ using the other variables can be removed by substituting $W$ and remembered for the inverse operation. In a crossing, start at either end of the underpass; the overpass keeps its label. At a curl with entry $a$ the new under label is $S_a^{\pm1}(a)=a$. If the returning overpass instead carries the unknown $b$, the equation $b=S_b^{\pm1}(a)$ gives $a=S_b^{\mp1}(b)=b$. In a bigon the lower strand is acted on by an operation and its inverse. In the triangle first name the top strand $z$, then propagate along the middle strand, then along the bottom one. Transport across the top sends the middle label $y$ to $A(y)$ with $A=S_z$ or $S_z^{-1}$. The two bottom computations agree because $A S_y^e=S_{A(y)}^e A$ for $e=1,-1$. The negative versions follow by inversion or substitution of $S_z^{-1}(y)$ in the given identity. This calculation can be read in either direction along each arc. All interior labels were solved for, so the two presentations reduce to identical boundary constraints in every orientation. Arbitrary outside relations are carried unchanged through these substitutions.
\end{proof}

\begin{theorem}[Wirtinger presentations]
Use one group generator for every arc between undercrossings, or one for a component with no undercrossings. At a crossing record the word relation $z=y^e x y^{-e}$ with $e=\pm1$ from the overpass normal. A move list gives explicit mutually inverse maps of the presented groups. The abelianization on a link of $r$ components is free abelian on its $r$ meridians. Mirroring preserves the isomorphism class of the group. Counting solutions in any supplied finite group table is a PR invariant.
\end{theorem}
\begin{proof}
One can equally put variables on all edge segments and impose equality across overpasses. In a group use $S_y(x)=yxy^{-1}$. Conjugation is invertible and preserves products and inverses; hence it commutes with conjugation by $y$ if the conjugating argument is transported to $S_z(y)$. Apply Lemma~\ref{lem:eliminate}. Each substitution yields words for images of generators, and substitution back gives the original words by a list of the same relation uses and free cancellations. Compose such lists along the movie. This is an entirely syntactic formulation: $u=v$ in a presentation means existence of a word-replacement certificate using the relators in contexts, free cancellations/insertions, and inverse words. Verification is bounded by the supplied code. In the abelianized presentation the crossing relation is simply $z=x$, so it identifies all arcs on a component. The map sending each arc to the corresponding standard basis vector in $\mathbb Z^r$ and its inverse choosing one representative per component give explicit inverse substitutions. For a mirror use the same overpasses with the ambient plane reflected, so the normal direction interchanges. Sending each generator to the inverse of its corresponding generator respects the reflected relations and squares to the identity. This reflected diagram gives the mirror as does switching all crossings (the two spatial reflections differ by an orientation-preserving half-turn). For a finite group of order $h$ test the relations on all $h^g$ assignments to $g$ generators. The local substitutions induce inverse bijections of solutions. The computation applies to any fixed union of conjugacy classes for the meridians as well, since propagation only conjugates them.
\end{proof}

For a trefoil a positive two-braid acts by $P(a,b)=(aba^{-1},a)$. The product $ab$ is fixed and $P^2$ conjugates both entries by $ab$, by cancellation. To close $P^3$ one needs its second entry $(ab)a(ab)^{-1}$ to equal $b$. Equality of the first then follows from the preserved product. Thus the trefoil has presentation
\[
 \langle a,b\mid aba=bab\rangle\ \simeq\ \langle u,v\mid u^2=v^3\rangle .
\]
Take $u=aba$, $v=ab$, so $u^2=(aba)(bab)=v^3$. Conversely $a=v^{-1}u$, $b=u^{-1}v^2$ give $ab=v$ and $aba=bab=u$. The transpositions $(12),(23)$ satisfy $aba=bab=(13)$ and generate the six permutations on three letters. This quotient is nonabelian, whereas every quotient of the one-generator unknot group is abelian. More explicitly, an alleged movie would convert this assignment to powers of a single permutation and then back, forcing its two entries to commute. No algorithm for equality of group words in an arbitrary presentation is involved. The classical identification with $\pi_1$ of the complement is the external Wirtinger theorem.

\subsection{An invariant from polynomial relations}
Put $R=\mathbb Z[t,t^{-1}]$. Use a generator for every segment and impose the two linear relations at a crossing: the over labels agree, and the under labels satisfy $z=t^e x+(1-t^e)y$. Give an isolated circle one free generator. The resulting row matrix $M$ has $g$ columns. We will use the greatest common divisor of its minors of size $g-1$, with the empty minor $1$. This is the diagrammatic Alexander polynomial up to multiplication by $\pm t^j$. The unreduced meridional coloring presentation has a free constant-color direction; the classical identification with covering-space homology concerns its reduced Alexander data.

\begin{lemma}[Minor calculus]
For any $j\ge0$, the common divisors of the minors of size $g-j$ are invariant under elementary row or column operations, insertion/deletion of a redundant relation, and adjoining a generator with a defining equation $v=W$. Minors of size at most zero generate the unit ideal; if the size exceeds the number of rows they generate the zero ideal. Indeed the corresponding ideals of minors agree.
\end{lemma}
\begin{proof}
A determinant is the alternating sum over permutations. Linearity in one row or column and cancellation of pairs when two rows agree follow directly by interchanging the selected entries. Thus adding a multiple of another row changes any minor by a multiple of another minor of the same size or by zero. The inverse operation gives the other containment. The same argument covers columns, a unit scaling, and adding a row already expressed as a linear combination of existing rows. To adjoin $v=W$, clear its defining row and column by these operations to get $\mathrm{diag}(M,1)$ with one more column. Its minors of the new size $k+1$, $k=g-j$, that use the last row and column are old $k$-minors. Those using exactly one vanish. Those using neither are $(k+1)$-minors of $M$, each a linear combination of $k$-minors by expansion. The size-zero convention covers $k\le0$. These arguments provide combinations witnessing both containments.
\end{proof}

\begin{theorem}[Alexander invariant and its specialization]
The first-minor gcd $\Delta_D$, taken up to $\pm t^j$, is invariant under oriented Reidemeister moves. For a knot it has a unique representative with least exponent zero and $\Delta_D(1)=1$. Mirroring sends the Laurent class to the class of $\Delta_D(t^{-1})$. A split diagram with two nonempty separated parts has zero first-minor gcd.
\end{theorem}
\begin{proof}
The affine actions $S_y(x)=tx+(1-t)y$ satisfy Lemma~\ref{lem:eliminate}: any affine map $A$ on the line gives $A S_y=S_{A(y)} A$. The substitutions and their inverses correspond precisely to adjoining or deleting defining equations. When the inner variables have been removed, the remaining crossing equations are equal as linear forms, up to rewriting by already imposed equalities and multiplication by units $t^e$. Such rewriting is elementary row addition or unit scaling. The minor calculus proves invariance, even of the whole family of determinantal ideals. A reflected planar drawing with the same overpass has the opposite normal, so $t$ is replaced by $t^{-1}$. In a split diagram each of the two blocks has row sum zero. Changing one column in each block to the sum of that block's columns produces two zero columns; hence every $(g-1)$-minor is zero.

For a knot, at $t=1$ the equations are exactly the differences of consecutive segment generators as one traverses the knot. For $g>1$ drop the final difference and one endpoint column. The remaining path matrix has determinant $\pm1$, by successive expansion at the path end. The $g=1$ case uses the empty minor. Thus there is a minor $m$ with $m(1)=\pm1$. The gcd divides it over the integral Laurent ring, so its value at $1$ is $\pm1$. Shifting the lowest exponent to zero and fixing its sign at $1$ gives the required normalization.
\end{proof}

\paragraph{A bounded gcd algorithm.}
Clear powers of $t$ separately in each nonzero minor so its least exponent is zero. Over $\mathbb Q[t]$, use long division and the Euclidean algorithm; degree drops at each nonzero remainder, so both loops use at most the starting degree plus two iterations per pair. Taking the primitive integral part of the last nonzero monic gcd and multiplying by the integer gcd of the contents of the nonzero inputs gives the integral gcd. If all minors vanish, return zero. For completeness, the product of primitive integer polynomials is primitive: if a prime $p$ divided every coefficient of the product, reducing modulo $p$ would give two nonzero factors with zero product, contrary to multiplication of their leading coefficients. A number $>1$ has a least divisor $>1$, prime by a bounded trial search. Consequently contents multiply, by factoring out the individual contents. Clearing rational denominators and comparing contents shows that a primitive divisor over $\mathbb Q[t]$ of an integral polynomial divides its primitive part over $\mathbb Z[t]$. Together with the two division identities for the Euclidean gcd, this proves the algorithm. The divisibility statement for any given candidate uses its code as a free parameter. The gcd has no factor $t$ since one shifted input already has nonzero constant term. The units of the Laurent ring are exactly $\pm t^j$ by the lowest and highest terms of a product.

For $g$ columns, minor enumeration uses bounded row/column masks and at most $g!$ permutations per determinant. Each matrix entry is a finite Laurent array, so a determinant has degree window at most $g$ times the input window after common clearing. Long division has a nested degree counter and ordinary rational coefficient operations, all PR. An arithmetic induction for polynomial division keeps $f=q h+r$ and a bound on the remaining degree; for Euclid it keeps the two successive division records (and optionally the two linear-combination multipliers) and the decreasing degree counter. The common-divisor argument takes a particular divisor as a free parameter, avoiding induction on an unbounded greatest-divisor property.

\subsection{The two-strand family and the figure-eight}
On a two-strand braid the action is now $P(x,y)=(x-t d,x)$ for $d=x-y$. Its new difference is $-td$. Set $S_m=\sum_{j=0}^{m-1}(-t)^j$. After $m$ positive crossings the change in the second entry is $S_m d$, and that in the first is $-t S_m d$. Eliminate the internal rows in the braid and use coordinates $(y,d)$. Closure imposes just $S_m d=0$, leaving $y$ free. The first minors have gcd $S_m$. In particular
\[
 \Delta_{T_{2,2k+1}}(t)=1-t+t^2-\cdots+t^{2k},\qquad
 |\Delta_{T_{2,2k+1}}(-1)|=2k+1 .
\]
This gives the determinant $3$ for the trefoil and $1$ for the unknot. At $t=-1$ the coloring relation recovers the Fox relation in Section~\ref{sec:fox}.

For the figure-eight, a block $Q=\sigma_1\sigma_2^{-1}$ sends $(x,y,z)$ to $((1-t)x+ty,z,t^{-1}x+(1-t^{-1})z)$. Put $y=x+d$, $z=x+e$. The first entry of $Q^2$ equals $x+t((1-t)d+e)$, so closure eliminates $e=(t-1)d$. The remaining two output-minus-input differences reduce to $t^{-1}(t^2-3t+1)d$ and $-t^{-2}(t^2-3t+1)d$. Thus the least-exponent-zero representative with value $1$ at $t=1$ is $-t^2+3t-1$. Its determinant is $5$. In particular the trefoil and figure-eight are inequivalent, as are their mirrors to each other across these two types.

\paragraph{The overview's Conway computation.}
For the odd two-strand knots the centered, value-one normalization is $t^{-k}S_{2k+1}$. For $k=1$ this is $t^{-1}-1+t=1+(t^{1/2}-t^{-1/2})^2$, giving the conventional Conway polynomial $1+z^2$ under $z=t^{1/2}-t^{-1/2}$. The same computation for the figure-eight gives $1-z^2$. Substitution is injective on $\mathbb Z[z]$, since the highest power in a nonzero polynomial contributes a unique top exponent in $t^{1/2}$. There is also a direct deduction of the entire two-strand recursion from the Conway axioms. Write $C_m$ for the Conway value of the standard closure, normalized by $C_1=1$. Applying the skein relation to a positive crossing gives $C_m=C_{m-2}+zC_{m-1}$ for $m\ge2$, since the switched pair cancels and the smoothing removes one crossing. The skein relation at a curl gives $zC_0=0$, hence $C_0=0$. Thus $C_2=z$, $C_3=1+z^2$, and $C_4=2z+z^3$. This final recurrence is a consequence of the stated Conway axioms; the diagrammatic Alexander-class construction and the explicit centered polynomials above were proved independently. We do not derive the globally sign-normalized Conway skein theorem for all links here.
